\section{Weekly Summary}

\begin{tcolorbox}[colback=yellow!10,colframe=black!50,title=AI-Generated Content]
\noindent The summary below was written by Claude (Anthropic) from that week's regression outputs and series descriptions. It has not been reviewed or fact-checked by a human, and should be read as an automated, informal recap -- not analysis.
\end{tcolorbox}

Welcome back to the weekly ritual where a computer flings two unrelated economic time series into a blender and I pretend to be surprised by what comes out. This week's contestants included Federal Reserve Treasury holdings, a discontinued Southern services price index, some population data from Nuckolls County, Nebraska (population: fewer than you'd hope), and a price index for congenital anomaly medical expenditures. Truly the crossover episodes nobody asked for. Let me pour one out for our dignity before we begin.

Monday kicked off with the Fed's medium-maturity Treasury holdings versus a CPI series so discontinued it retired back in 2017. The raw regression came out swaggering with an R-squared of 0.77, as if it had discovered something. Then we detrended it, the R-squared slumped to 0.37, and the coefficient literally flipped sign from positive to negative. Nothing says "we were just watching two things drift upward together" quite like a relationship that can't even decide which direction it points once you take away its little time crutch. Classic trend mirage. Dunk registered. Tuesday, however, actually behaved: nonfinancial noncorporate business quasi-corporation withdrawals versus current-cost depreciation of nonresidential structures held onto an R-squared of 0.87 even after detrending, and the coefficient barely budged from 0.42 to 0.40. That's the rare, annoying case where the relationship survives the trend-stripping, probably because both are just proxies for "big dollar amounts sloshing around the business economy." I'm reluctantly impressed, which is a feeling I try to avoid on this project.

Wednesday gave us Nuckolls County's dwindling population predicting a congenital anomaly price index, and the raw fit was a dramatic 0.88 — a coefficient of negative 75, meaning as Nebraskans leave, medical prices climb, which is exactly the kind of ludicrous causal story I'm legally required to tell you is nonsense. Detrended, it slumped to 0.36 but technically stayed "significant," so it's a mirage wearing a slightly convincing hat. Thursday was refreshingly humble: net unilateral transfers versus bank failure transactions produced a raw R-squared of 0.005, which is the statistical equivalent of a shrug, though the detrended version squeaked out a p-value just under 0.05, a "significance" so thin you could read a newspaper through it. And Friday closed strong with government investment in intellectual property versus manufacturing compensation costs — 0.96 raw, 0.84 detrended, coefficient holding steady. Another one that stubbornly refuses to be dismissed.

As always, the mandatory reminder: these are randomly paired series with zero causal intent and no rigorous statistical claim whatsoever. When a relationship "holds up" here, it means the two things happened to wiggle together, not that one causes the other, and not that you should build a portfolio, a policy, or a personality around it. Spurious correlation is the house special here, served daily. This week just happened to plate a couple of dishes that looked suspiciously edible. Don't eat them anyway. See you next week for more distinguished nonsense.
